\section{Estructura de la memoria}

La presente memoria se organiza en seis capítulos y tres anexos. El presente capítulo
recoge la introducción y motivación del trabajo, los objetivos específicos y la
estructura de la memoria.

El Capítulo 2 presenta el estado del arte en producción audiovisual: la
evolución desde la producción tradicional hasta el modelo remoto basado en \textit{software},
los protocolos de red y los estándares de contribución y distribución de vídeo, y una
comparativa de enfoques y herramientas de producción remota por \textit{software}. Se incluyen también casos de uso reales de
producción remota a escala industrial.

El Capítulo 3 describe el diseño e implementación del sistema Voctomix~2.0:
la arquitectura cliente-servidor (\texttt{voctocore} y \texttt{voctogui}), las fuentes de señal de
cámaras y auxiliares, los modos de composición, las transiciones, el sistema de control
de estados de emisión (\textit{stream blanker}), el sistema de \textit{overlays} y rotulación
dinámica, y el servicio de telemetría mediante RabbitMQ.

El Capítulo 4 detalla el despliegue del sistema en sus distintos escenarios: en entorno nativo tanto en un único equipo como en dos equipos, el despliegue contenerizado con Docker~Compose
y el despliegue orquestado con Kubernetes.

El Capítulo 5 recoge las pruebas de validación del sistema en los escenarios
de despliegue, con mediciones de latencia de conmutación, consumo de recursos, estabilidad
a largo plazo, resiliencia ante fallos y registros de telemetría.

El Capítulo 6 presenta las conclusiones del trabajo y las líneas de trabajo
futuro identificadas.

Los Anexos incluyen los aspectos éticos, económicos, sociales y ambientales
del proyecto (Anexo~A), el presupuesto económico detallado (Anexo~B)
y la evolución del repositorio a lo largo del desarrollo del proyecto (Anexo~C).